\documentclass[10pt,a4paper]{article}
\usepackage[T1]{fontenc}
\usepackage[utf8]{inputenc}
\usepackage{newtxtext,newtxmath}
\usepackage[a4paper,margin=18mm]{geometry}
\usepackage{amsmath}
\usepackage{graphicx}
\usepackage{microtype}
\usepackage{booktabs,threeparttable,tabularx,array,longtable}
\usepackage{float}
\usepackage{needspace}
\usepackage[font=small,labelfont=bf,justification=raggedright,singlelinecheck=false]{caption}
\usepackage[superscript]{cite}
\usepackage[hidelinks]{hyperref}
\hypersetup{pdftitle={Modelled prevalence of colorectal adenomas and cancer in PRENEC},pdfauthor={PRENEC Study Group}}
\setlength{\parindent}{0pt}
\setlength{\parskip}{5pt}
\setlength{\emergencystretch}{2em}
\setcounter{secnumdepth}{0}
\renewcommand{\arraystretch}{1.12}
\newcommand{\Se}{\mathit{Se}}
\newcommand{\Sp}{\mathit{Sp}}
\newcommand{\Adh}{\operatorname{Adh}}
\newcommand{\BetaDist}{\operatorname{Beta}}
\newcommand{\dec}{\ensuremath{\cdot}}
\newcolumntype{C}{>{\centering\arraybackslash}p{0.17\linewidth}}
\newcolumntype{Y}{>{\raggedright\arraybackslash}X}
\newcolumntype{L}[1]{>{\raggedright\arraybackslash}p{#1}}
\makeatletter
\setlength{\@fptop}{0pt}
\makeatother

\title{\textbf{Modelled prevalence of colorectal adenomas and cancer in a public-sector screening programme in Chile: a retrospective analysis of PRENEC}}
\author{PRENEC Study Group}
\date{}

\begin{document}
\maketitle

\section*{Summary}

\textbf{Background}\quad
PRENEC used FIT to select participants for colonoscopy, so observed lesion
counts did not directly estimate prevalence. This study estimated patient-level
prevalence and detection expected under universal colonoscopy for low-risk
adenoma, high-risk adenoma, and colorectal cancer.

\textbf{Methods}\quad
Routinely collected records were analysed for public-insurance beneficiaries
aged 50--75 years enrolled in eight regions during 2012--2022. Participants
without a reported family history of colorectal cancer formed the target
population. One available OC-Sensor faecal immunochemical test (FIT) was
randomly selected per person and classified as positive at $\geq100$ ng Hb/mL.
Observed detection was corrected for lesion-specific FIT sensitivity and
colonoscopy participation; prevalence was additionally corrected for
colonoscopy sensitivity. PRENEC-derived and literature-based FIT sensitivities
defined two scenarios. Uncertainty was propagated through 10\,000 Monte Carlo
iterations.

\textbf{Findings}\quad
Among 34\,241 participants, 25\,014 (73\dec{}1\%) were women. PRENEC-derived
FIT sensitivity was 7\dec{}0\% for low-risk adenoma,
21\dec{}8\% for high-risk adenoma, and 100\dec{}0\% for colorectal cancer.
In assignment A, adjusted population detection was 20\dec{}9\%
(95\% HPD interval 9\dec{}0--35\dec{}1), 9\dec{}0\%
(4\dec{}9--13\dec{}7), and 0\dec{}271\% (0\dec{}040--0\dec{}558),
respectively. Corresponding modelled prevalence was 26\dec{}1\%
(11\dec{}5--44\dec{}2), 11\dec{}1\% (5\dec{}7--17\dec{}1), and
0\dec{}300\% (0\dec{}049--0\dec{}623). Literature-based estimates were
23\dec{}6\%, 10\dec{}4\%, and 0\dec{}408\%; assignment B was similar.

\textbf{Interpretation}\quad
Modelled adenoma prevalence was substantial, whereas colorectal cancer
prevalence was below 1\%. Estimates were stable across FIT-sensitivity
scenarios, but uncertainty and transportability assumptions remain important.

\textbf{Funding}\quad
Funding information is reported separately in the submission metadata.

\section*{Research in context}

\textbf{Evidence before this study}\quad
Colorectal lesion prevalence has been estimated from autopsy series and
cross-sectional colonoscopy studies, but estimates vary with participant
selection, lesion definitions, and incomplete detection at colonoscopy.
Published Latin American evidence has mainly described adenoma detection
rather than underlying prevalence. Chilean evidence has been limited to a
small primary-colonoscopy study and partial reports from PRENEC.

\textbf{Added value of this study}\quad
This study used a large public-sector screening cohort from eight Chilean
regions and explicitly separated three quantities: lesion yield among
colonoscopy completers, population detection after correction for FIT
sensitivity and colonoscopy participation, and prevalence after additional
correction for colonoscopy sensitivity. PRENEC-derived and external FIT
sensitivity inputs were compared, and all principal sources of parameter
uncertainty, including colonoscopy participation, were propagated.

\Needspace{7\baselineskip}
\textbf{Implications of all the available evidence}\quad
Local model-based estimates can inform the expected yield and colonoscopy
requirements of colorectal cancer screening in Chile. Their use in policy or
cost-effectiveness models should retain uncertainty around FIT and
colonoscopy sensitivity and should not assume national representativeness
without additional standardisation.

\section*{Introduction}

Colorectal cancer incidence and mortality vary markedly across countries,
with Chile historically reporting lower rates than many high-income
settings.\cite{Arnold2017,IARC2026} These differences raise questions about
the applicability of international estimates of adenoma prevalence to the
Chilean population. Cancer rates reflect both the occurrence of precursor
lesions and their progression, and cannot establish adenoma prevalence
directly. Local estimates are therefore needed to inform models of colorectal
cancer natural history and screening policy. These parameters are important
for predicting screening yield, planning colonoscopy provision, and
evaluating the cost-effectiveness of alternative screening strategies.

The prevalence of colorectal precursor lesions and cancer has been
investigated through autopsy studies and cross-sectional studies of
participants undergoing colonoscopy.\cite{Wang2023,Heitman2009} Estimates
vary with age, participant selection, lesion definitions, and ascertainment
methods. Colonoscopy studies also underestimate true prevalence when lesions
remain undetected.\cite{Zhao2019} A systematic review and meta-analysis by
Montalvan-Sanchez and colleagues reported adenoma detection rates in Latin
America broadly comparable to those in high-income countries: 19\dec{}9\%
with primary screening colonoscopy and 39\dec{}0\% after a positive
FIT.\cite{Montalvan2024}

In Chile, no autopsy-based estimates of colorectal lesion prevalence have
been identified, whereas primary screening colonoscopy has been studied only
in a small private-clinic sample in Santiago, limiting inference about the
general adult population.\cite{Silva2011}

PRENEC was a large public-sector colorectal cancer screening pilot developed
through collaboration between Chile and Japan.\cite{Okada2016} It enrolled
asymptomatic beneficiaries of the public health insurance system (FONASA)
aged 50--75 years across eight regions during 2012--2022. Existing reports
described programme implementation and screening performance but did not
estimate underlying lesion prevalence or population detection under
universal colonoscopy. Because the programme used FIT to select most
participants for colonoscopy, its observed lesion counts required explicit
modelling of incomplete ascertainment.

This study aimed to estimate the prevalence of low-risk adenoma, high-risk
adenoma, and colorectal cancer, and the population detection expected under
universal colonoscopy, among PRENEC participants without a family history of
colorectal cancer. A probabilistic framework linked observed detection to
FIT sensitivity, colonoscopy participation, and colonoscopy sensitivity.
Analyses used PRENEC-derived and literature-derived FIT sensitivities and
were conducted for both sexes combined and separately for women and men.

\section*{Methods}

\subsection*{Study design and participants}

This retrospective analysis used routinely collected records from PRENEC, a
multicentre colorectal cancer screening programme operating in eight Chilean
regions. PRENEC enrolled asymptomatic FONASA beneficiaries aged 50--75 years
between 2012 and 2022; records of first colonoscopies were available through
March, 2023. The estimands pertained to the eligible PRENEC population and
were not interpreted as nationally representative.

Participants without a reported family history of colorectal cancer formed
the target screening population. Participants with a family history were
excluded from the target prevalence denominator but provided an internal
cohort for estimating FIT sensitivity because colonoscopy was offered
irrespective of FIT results. The target-population denominator, $N$, included
all eligible participants aged 50--75 years without a family history,
irrespective of whether a FIT result was recorded. No additional correction
for FIT coverage was applied.

PRENEC supplied two OC-Sensor FIT kits. Under the programme protocol,
colonoscopy was offered if either result was at least 100 ng Hb/mL. The
descriptive programme analysis therefore used positivity in either round. For
the modelled single-sample strategy, one available FIT result was randomly
selected per participant. Participants without an available FIT remained in
$N$ but could not enter the detected pathway.

Two complementary FIT assignments, A and B, were generated using the
prespecified seed 20260923. When two results were available, one was allocated
to A at random and the other to B; when only one result was available, both
assignments used it. The assignments contained the same participants and
were not independent samples. Assignment A was prespecified for narrative
reporting and assignment B was retained as a robustness analysis.

A qualifying colonoscopy occurred from the first colonoscopy-check date
(inclusive) to the second colonoscopy-check date (exclusive). When the second
date was missing, no upper boundary was applied. The first qualifying
colonoscopy was used.

These eligibility, FIT-selection, and colonoscopy-window definitions specified
the observed detection pathway. The observed population detection rate was
then corrected for FIT sensitivity and colonoscopy participation to estimate
adjusted population detection, and for colonoscopy sensitivity to estimate
modelled lesion prevalence. The following sections define the lesion categories
and target quantities, specify the model inputs, and describe uncertainty
propagation.

\subsection*{Lesion ascertainment and outcomes}

Colonoscopic findings were described with the Paris classification, and
histopathological findings were classified according to the Vienna
classification and the standardised PRENEC pathology
protocol.\cite{Paris2003,Vienna2000,Ito2016,Okada2016} Analyses were
person-level and required colonoscopic detection with histological
confirmation. When several lesions were identified, PRENEC retained the
characteristics of the lesion with the highest recorded risk.

The publication analysis considered low-risk adenoma, high-risk adenoma, and
colorectal cancer. Low-risk adenomas were tubular adenomas with low-grade
dysplasia smaller than 10 mm. High-risk lesions comprised tubular adenomas
with low-grade dysplasia measuring at least 10 mm; tubular adenomas with
high-grade dysplasia; tubulovillous adenomas with low-grade or high-grade
dysplasia; villous adenomas with high-grade dysplasia; and intramucosal
carcinoma. Colorectal cancer comprised well, moderately, or poorly
differentiated invasive adenocarcinoma. Adenoma number did not enter the risk
classification. External advanced-adenoma evidence was used only as an
explicit proxy where stated.

\subsection*{Target quantities and detection pathway}

For lesion category $j$, $p_j$ denoted the probability that a participant in
the target population had at least one true lesion of category $j$. Let
$D_j$ denote the observed number of participants whose selected FIT was
positive, who completed colonoscopy in the prespecified window, and in whom
at least one lesion of category $j$ was histologically confirmed. The
observed programme detection rate for the single-FIT pathway was

\begin{equation}
d_j=\frac{D_j}{N}.
\label{eq:d}
\end{equation}

The sequential pathway was represented as

\begin{equation}
d_j=\frac{D_j}{N}
=p_j\,\Se_{\mathrm{FIT},j}\,
\bigl(\Adh\mid\Se_{\mathrm{FIT},j}\bigr)\,
\bigl(\Se_{\mathrm{COL},j}\mid\Adh,\Se_{\mathrm{FIT},j}\bigr),
\label{eq:conditional}
\end{equation}

where $\Se_{\mathrm{FIT},j}$ was the person-level sensitivity of the selected
FIT for lesion category $j$, $\Adh$ was colonoscopy participation conditional
on a positive selected FIT, and $\Se_{\mathrm{COL},j}$ was person-level
colonoscopy sensitivity for detecting at least one lesion of category $j$.
This factorisation described successive conditional stages; it did not
require unconditional independence between them.\cite{Fanshawe2024}

The model assumed that colonoscopy participation did not vary with FIT
sensitivity and that colonoscopy sensitivity did not vary with either FIT
sensitivity or colonoscopy participation:

\begin{align}
\bigl(\Adh\mid\Se_{\mathrm{FIT},j}\bigr)&=\Adh,\\
\bigl(\Se_{\mathrm{COL},j}\mid\Adh,\Se_{\mathrm{FIT},j}\bigr)
&=\Se_{\mathrm{COL},j}.
\end{align}

The plausibility of these assumptions and the robustness of the estimates to
potential departures from them are considered in the Discussion.

Within sex and age strata, observed colonoscopy participation was assumed to
represent participation among lesion carriers. FIT sensitivity from the
internal cohort or literature and colonoscopy sensitivity from external
evidence were assumed transportable to the target population. Lesion
definitions and units of analysis were assumed sufficiently comparable.
Colonoscopy plus histological ascertainment was treated as effectively
specific. Under these assumptions,

\begin{equation}
d_j=p_j\,\Se_{\mathrm{FIT},j}\,\Adh\,\Se_{\mathrm{COL},j}.
\label{eq:pathway}
\end{equation}

Figure~\ref{fig:pathway} summarises the observed and counterfactual pathways,
the target quantities, and the conditional-independence assumptions.

\begin{figure}[p]
\centering
\includegraphics[width=\textwidth,height=0.78\textheight,keepaspectratio]
  {output/pdf/prenec_probability_tree.pdf}
\caption{\textbf{Person-level lesion-detection pathway for category $j$.}
The solid path represents observed detection through the selected FIT,
colonoscopy participation, and colonoscopy sensitivity. The dashed path
represents the counterfactual expected count under universal colonoscopy.
Lower-case symbols denote rates or probabilities and upper-case symbols
denote counts. Hats are omitted throughout the diagram. FIT=faecal
immunochemical test; $\Adh$=colonoscopy participation after a positive
selected FIT.}
\label{fig:pathway}
\end{figure}

\subsection*{Adjusted population detection and prevalence}

Let $R_j$ denote the expected number of participants in whom lesion category
$j$ would be detected if all members of the target population underwent
colonoscopy of the quality represented by the model. The adjusted population
detection rate was

\begin{equation}
r_j=\frac{R_j}{N}
=p_j\,\Se_{\mathrm{COL},j}
=\frac{d_j}{\Se_{\mathrm{FIT},j}\,\Adh}.
\label{eq:r}
\end{equation}

This quantity corrects for FIT sensitivity and colonoscopy participation but
retains lesions missed by colonoscopy. Modelled lesion prevalence was

\begin{equation}
p_j=
\frac{d_j}
{\Se_{\mathrm{FIT},j}\,\Adh\,\Se_{\mathrm{COL},j}}
=\frac{r_j}{\Se_{\mathrm{COL},j}}.
\label{eq:p}
\end{equation}

The observed colonoscopy detection rate in the selected-FIT pathway was

\begin{equation}
\mathrm{cdr}_j=\frac{D_j}{C},
\end{equation}

where $C$ was the number of selected-FIT-positive participants completing
colonoscopy. This descriptive rate received no sensitivity correction.
The relation of equation~\eqref{eq:p} to the Rogan--Gladen correction is
shown in the Supplementary Methods.\cite{RoganGladen1978}

\begin{figure}[p]
\centering
\includegraphics[width=\textwidth,height=0.84\textheight,keepaspectratio]
  {output/pdf/prenec_methodology_framework.pdf}
\caption{\textbf{PRENEC information sources and population estimators.}
The target population provided $D_j$, $d_j$, and $\Adh$; the
family-history cohort provided PRENEC-derived FIT sensitivity. External
evidence supplied an alternative FIT-sensitivity scenario and
lesion-specific colonoscopy sensitivity. The framework was applied
separately to low-risk adenoma, high-risk adenoma, and colorectal cancer, for
both sexes combined, women, and men.}
\label{fig:framework}
\end{figure}

Figure~\ref{fig:framework} summarises the two analytic populations and the
source of each empirical or external model input.

\subsection*{Model inputs}

\textit{FIT sensitivity.}\quad
PRENEC-derived sensitivity was estimated in participants with a family
history who had an evaluable FIT and colonoscopy. For each lesion category,
sensitivity was the proportion with a positive randomly selected FIT among
participants classified in that category at colonoscopy. Wilson intervals
were used descriptively. The external low-risk adenoma estimate came from a
four-study random-effects meta-analysis of OC-Sensor data at the programme
threshold.\cite{Park2010,Redwood2016,Chang2017,Imperiale2014} Published
OC-Sensor estimates from Lin and colleagues were used for colorectal cancer
and, with advanced adenoma as an explicit proxy, for high-risk
adenoma.\cite{Lin2021}

\textit{Colonoscopy sensitivity.}\quad
External diagnostic-accuracy and tandem-colonoscopy evidence informed
colonoscopy sensitivity.\cite{Martin2013,vanRijn2006,Zhao2019,Zhao2023}
The low-risk input was a composite proxy combining size and histology
evidence; the high-risk input used the patient-level advanced-adenoma result
from Zhao and colleagues as an explicit proxy for the PRENEC high-risk
category; and the colorectal cancer input used 52 detected cancers among 57
patients from the Martin-Lopez/AETSA synthesis. These parameters were used
only in equation~\eqref{eq:p}, not in $\mathrm{cdr}_j$ or $r_j$.

\textit{Colonoscopy participation.}\quad
For each complementary assignment, sex, and age group, participation was

\begin{equation}
\Adh=\frac{C}{F^+},
\end{equation}

where $F^+$ was the number positive on the selected FIT and $C$ the number
completing colonoscopy in the prespecified window. Age groups were 50--54,
55--59, 60--64, 65--69, 70--74, and 75 years.

\subsection*{Statistical analysis}

All quantities were estimated for both sexes combined and separately for
women and men. Age-specific lesion-detection draws were corrected with the
corresponding age-group colonoscopy-participation draws and then weighted by
the age distribution of the target population. Results were reported for
assignment A, with assignment B displayed as a paired robustness analysis.

Uncertainty was propagated through 10\,000 Monte Carlo iterations. Within age
$e$, observed detection was represented by
$\BetaDist(D_{j,e},N_e-D_{j,e})$ without pseudocounts. A shared uniform
quantile was applied across ages within each lesion and iteration, preserving
the implemented rank dependence. PRENEC-derived FIT sensitivity followed
$\BetaDist(\mathrm{TP}_j+0\dec5,\mathrm{FN}_j+0\dec5)$, except that the
PRENEC cancer sensitivity was fixed at 1\dec0. Colonoscopy participation was
sampled from
$\BetaDist(C_e+0\dec5,F^+_e-C_e+0\dec5)$ for each age group; the same
participation quantiles were reused across lesion categories and FIT
sensitivity scenarios. Colonoscopy sensitivities followed the beta
distributions in the parameter workbook. Published pooled FIT estimates used
moment-matched beta distributions, and the de novo low-risk meta-analysis
used a Hartung--Knapp-calibrated logit-normal distribution.

Simulated proportions were not truncated at 1. Point estimates were
simulation means and uncertainty was summarised with 95\% highest posterior
density (HPD) intervals. Seeds were prespecified for FIT selection, model
simulation, colonoscopy sensitivity, external FIT sensitivity, and
colonoscopy participation. Analyses were performed in R and structured for
reporting with STROBE and RECORD; parameter provenance followed relevant
GATHER principles.\cite{STROBE2007,RECORD2015,GATHER2016}

\clearpage
\section*{Results}

\subsection*{Study population and screening outcomes}

The target population comprised 34\,241 PRENEC participants aged 50--75
years without a reported family history of colorectal cancer; 25\,014
(73\dec{}1\%) were women and mean age was 60\dec{}8 years (95\% CI
60\dec{}7--60\dec{}8; table~\ref{tab:screening}). Both FIT rounds were
recorded for 31\,095 (90\dec{}8\%) participants. A concentration at or above
100 ng Hb/mL was recorded in 8\dec{}8\% in round 1 and 8\dec{}9\% in round
2; concentrations were higher among men.

At least one programme FIT was positive for 4\,087 participants, of whom
3\,433 (84\dec{}0\%) completed a qualifying colonoscopy. There were 3\,434
qualifying colonoscopies in total, including one in a participant without a
positive result in either recorded round. Adequate bowel preparation was
recorded for 79\dec{}4\%, caecal intubation for 95\dec{}0\%, and a withdrawal
time of at least 8 min for 90\dec{}2\%. In the descriptive programme pathway,
low-risk adenoma was found in 18\dec{}1\%, high-risk adenoma in 22\dec{}1\%,
and colorectal cancer in 2\dec{}4\% of qualifying colonoscopies. Detection
was numerically higher among men.

\begin{table}[p]
\centering
\caption{Screening characteristics and colonoscopy outcomes among PRENEC
participants aged 50--75 years without a reported family history of
colorectal cancer}
\label{tab:screening}
\begin{threeparttable}
\scriptsize
\setlength{\tabcolsep}{3pt}
\renewcommand{\arraystretch}{1.04}
\begin{tabularx}{\linewidth}{@{}>{\raggedright\arraybackslash}X C C C@{}}
\toprule
 & \textbf{Both sexes} & \textbf{Women} & \textbf{Men}\\
\midrule
\textit{\textbf{Participants}} & & &\\
Participants, n & 34\,241 & 25\,014 & 9\,227\\
Age, years, mean (95\% CI) & 60\dec{}8 (60\dec{}7--60\dec{}8) &
60\dec{}5 (60\dec{}4--60\dec{}6) &
61\dec{}5 (61\dec{}4--61\dec{}7)\\
\addlinespace
\textit{\textbf{FIT}} & & &\\
FIT $\geq100$ ng Hb/mL, round 1, n/N (\%) &
2\,727/31\,102 (8\dec{}8\%) &
1\,844/22\,835 (8\dec{}1\%) &
883/8\,267 (10\dec{}7\%)\\
FIT $\geq100$ ng Hb/mL, round 2, n/N (\%) &
2\,755/31\,097 (8\dec{}9\%) &
1\,838/22\,832 (8\dec{}1\%) &
917/8\,265 (11\dec{}1\%)\\
FIT concentration, round 1, ng Hb/mL, mean (95\% CI) &
84\dec{}9 (79\dec{}7--90\dec{}2) &
74\dec{}6 (68\dec{}9--80\dec{}4) &
113\dec{}5 (101\dec{}5--125\dec{}5)\\
FIT concentration, round 2, ng Hb/mL, mean (95\% CI) &
89\dec{}8 (84\dec{}3--95\dec{}4) &
79\dec{}8 (73\dec{}6--86\dec{}0) &
117\dec{}6 (105\dec{}4--129\dec{}8)\\
Both FIT rounds recorded, n/N (\%) &
31\,095/34\,241 (90\dec{}8\%) &
22\,830/25\,014 (91\dec{}3\%) &
8\,265/9\,227 (89\dec{}6\%)\\
\addlinespace
\textit{\textbf{Colonoscopy participation}} & & &\\
Participation after any positive FIT, n/N (\%) &
3\,433/4\,087 (84\dec{}0\%) &
2\,366/2\,794 (84\dec{}7\%) &
1\,067/1\,293 (82\dec{}5\%)\\
\addlinespace
\textit{\textbf{Colonoscopy quality}} & & &\\
Boston bowel-preparation score 8--9, n/N (\%) &
2\,727/3\,434 (79\dec{}4\%) &
1\,890/2\,367 (79\dec{}8\%) &
837/1\,067 (78\dec{}4\%)\\
Caecal intubation, n/N (\%) &
3\,261/3\,434 (95\dec{}0\%) &
2\,246/2\,367 (94\dec{}9\%) &
1\,015/1\,067 (95\dec{}1\%)\\
Withdrawal time $\geq8$ min, n/N (\%) &
3\,099/3\,434 (90\dec{}2\%) &
2\,131/2\,367 (90\dec{}0\%) &
968/1\,067 (90\dec{}7\%)\\
\addlinespace
\textit{\textbf{Programme colonoscopy yield}} & & &\\
Low-risk adenoma, n/N (\%) &
623/3\,434 (18\dec{}1\%) &
408/2\,367 (17\dec{}2\%) &
215/1\,067 (20\dec{}1\%)\\
High-risk adenoma, n/N (\%) &
759/3\,434 (22\dec{}1\%) &
476/2\,367 (20\dec{}1\%) &
283/1\,067 (26\dec{}5\%)\\
Colorectal cancer, n/N (\%) &
83/3\,434 (2\dec{}4\%) &
43/2\,367 (1\dec{}8\%) &
40/1\,067 (3\dec{}7\%)\\
\bottomrule
\end{tabularx}
\begin{tablenotes}[flushleft]
\scriptsize
\item[] Data are n/N (\%) or mean (95\% CI). FIT-round denominators are
participants with a recorded result in that round. Both-round completion uses
the full target population.
\item[] Colonoscopy participation used participants with at least one FIT
$\geq100$ ng Hb/mL. Quality and lesion yield used all 3\,434 qualifying first
colonoscopies; one occurred outside the programme-positive subgroup.
\item[] Quality-indicator denominators were participants with a recorded
value. The source withdrawal-time field is labelled $>8$ min but records at
exactly 8 min are coded as meeting the criterion; the observed coding is
therefore reported as $\geq8$ min.
\item[] FIT=faecal immunochemical test; CI=confidence interval.
\end{tablenotes}
\end{threeparttable}
\end{table}

\subsection*{Sensitivity estimates}

PRENEC-derived FIT sensitivity was 7\dec{}0\% (95\% CI
4\dec{}6--10\dec{}4) for low-risk adenoma and 21\dec{}8\%
(16\dec{}2--28\dec{}6) for high-risk adenoma
(table~\ref{tab:localfit}). All six participants classified with colorectal
cancer in the internal cohort had a positive selected FIT, yielding a point
estimate of 100\dec{}0\% with a wide Wilson interval
(61\dec{}0--100\dec{}0).

External FIT sensitivities were similar for both adenoma categories:
7\dec{}6\% for low-risk adenoma and 23\dec{}0\% for high-risk adenoma
(table~\ref{tab:external}). For colorectal cancer, the external estimate was
74\dec{}0\%, lower than the PRENEC point estimate but within its wide
confidence interval. External colonoscopy sensitivity ranged from
79\dec{}8\% to 91\dec{}2\%; the lesion-specific sources included explicit
proxies and indirect evidence.

\begin{table}[p]
\centering
\caption{Patient-level FIT sensitivity estimated in the PRENEC
family-history cohort}
\label{tab:localfit}
\begin{threeparttable}
\small
\setlength{\tabcolsep}{7pt}
\renewcommand{\arraystretch}{1.15}
\begin{tabularx}{\linewidth}{@{}Y c c c@{}}
\toprule
\textbf{Lesion category} & \textbf{FIT-positive cases, n} &
\textbf{Cases, N} & \textbf{Sensitivity, \% (95\% CI)}\\
\midrule
Low-risk adenoma & 21 & 301 & 7\dec{}0 (4\dec{}6--10\dec{}4)\\
High-risk adenoma & 37 & 170 & 21\dec{}8 (16\dec{}2--28\dec{}6)\\
Colorectal cancer & 6 & 6 & 100\dec{}0 (61\dec{}0--100\dec{}0)\\
\bottomrule
\end{tabularx}
\begin{tablenotes}[flushleft]
\footnotesize
\item[] One observed FIT measurement was randomly selected per participant
and classified at $\geq100$ ng Hb/mL. Confidence intervals were calculated
with the Wilson method. The fixed FIT-selection seed was 20260923.
\item[] FIT=faecal immunochemical test; CI=confidence interval.
\end{tablenotes}
\end{threeparttable}
\end{table}

\begin{table}[p]
\centering
\caption{External FIT and colonoscopy sensitivity inputs used in the
prevalence models}
\label{tab:external}
\begin{threeparttable}
\scriptsize
\setlength{\tabcolsep}{4pt}
\renewcommand{\arraystretch}{1.13}
\begin{tabularx}{\linewidth}{@{}L{0.18\linewidth}L{0.18\linewidth}
L{0.20\linewidth}L{0.18\linewidth}Y@{}}
\toprule
\textbf{Lesion category} &
\textbf{External FIT sensitivity, \% (95\% CI)} &
\textbf{FIT evidence} &
\textbf{External colonoscopy sensitivity, \% (95\% CI)} &
\textbf{Colonoscopy evidence}\\
\midrule
Low-risk adenoma &
7\dec{}6 (6\dec{}4--8\dec{}9) &
Four-study OC-Sensor meta-analysis &
79\dec{}8 (76\dec{}8--82\dec{}6) &
Composite size-histology proxy\\
High-risk adenoma &
23\dec{}0 (20\dec{}0--25\dec{}0) &
Lin et al., 2021; advanced adenoma proxy &
82\dec{}1 (68\dec{}0--91\dec{}6) &
Zhao et al., 2023; advanced adenoma proxy (32/39)\\
Colorectal cancer &
74\dec{}0 (64\dec{}0--83\dec{}0) &
Lin et al., 2021 &
91\dec{}2 (81\dec{}8--96\dec{}6) &
Martin-Lopez/AETSA, 2013 (52/57)\\
\bottomrule
\end{tabularx}
\begin{tablenotes}[flushleft]
\footnotesize
\item[] The low-risk FIT estimate came from a random-effects meta-analysis at
an OC-Sensor threshold of 100 ng Hb/mL. Lin et al. reported OC-Sensor results
at 20 $\mu$g Hb/g faeces.
\item[] Colonoscopy inputs combined direct and modelled evidence. The
low-risk estimate was a composite proxy; the high-risk estimate used advanced
adenoma as a proxy; and the cancer estimate was based on a small indirect
sample.
\item[] FIT=faecal immunochemical test; CI=confidence interval.
\end{tablenotes}
\end{threeparttable}
\end{table}

\subsection*{Observed detection, adjusted detection, and prevalence}

In assignment A, 2\,763 participants had a positive selected FIT and 2\,323
completed colonoscopy in the modelled window. Corresponding counts in
assignment B were 2\,760 and 2\,315. The uncorrected selected-pathway
colonoscopy detection rates in assignment A were 17\dec{}5\% for low-risk
adenoma, 23\dec{}9\% for high-risk adenoma, and 3\dec{}36\% for colorectal
cancer; assignment B estimates were 17\dec{}1\%, 25\dec{}0\%, and
3\dec{}15\%, respectively.

Using PRENEC-derived FIT sensitivity, adjusted population detection in
assignment A was 20\dec{}9\% (95\% HPD interval 9\dec{}0--35\dec{}1) for
low-risk adenoma, 9\dec{}0\% (4\dec{}9--13\dec{}7) for high-risk adenoma,
and 0\dec{}271\% (0\dec{}040--0\dec{}558) for colorectal cancer
(table~\ref{tab:modelled}; figure~\ref{fig:modelled}). After additional
correction for colonoscopy sensitivity, modelled prevalence was 26\dec{}1\%
(11\dec{}5--44\dec{}2), 11\dec{}1\% (5\dec{}7--17\dec{}1), and
0\dec{}300\% (0\dec{}049--0\dec{}623), respectively.

Literature-derived FIT sensitivity produced slightly lower estimates for both
adenoma categories and a higher estimate for colorectal cancer. Modelled
prevalence was 23\dec{}6\% for low-risk adenoma, 10\dec{}4\% for high-risk
adenoma, and 0\dec{}408\% for colorectal cancer. Differences between
complementary assignments were small relative to the uncertainty intervals:
for both sexes, the largest difference in mean modelled prevalence was
0\dec{}53 percentage points.

\begin{table}[p]
\centering
\caption{Adjusted population detection and modelled prevalence in assignment
A, both sexes}
\label{tab:modelled}
\begin{threeparttable}
\scriptsize
\setlength{\tabcolsep}{4pt}
\renewcommand{\arraystretch}{1.12}
\begin{tabularx}{\linewidth}{@{}Y L{0.19\linewidth}L{0.19\linewidth}
L{0.19\linewidth}L{0.19\linewidth}@{}}
\toprule
& \multicolumn{2}{c}{\textbf{Adjusted population detection, \%}} &
\multicolumn{2}{c}{\textbf{Modelled prevalence, \%}}\\
\cmidrule(lr){2-3}\cmidrule(lr){4-5}
\textbf{Lesion category} &
\textbf{PRENEC FIT sensitivity} &
\textbf{Literature FIT sensitivity} &
\textbf{PRENEC FIT sensitivity} &
\textbf{Literature FIT sensitivity}\\
\midrule
Low-risk adenoma &
20\dec{}9 (9\dec{}0--35\dec{}1) &
18\dec{}8 (10\dec{}0--29\dec{}2) &
26\dec{}1 (11\dec{}5--44\dec{}2) &
23\dec{}6 (12\dec{}3--36\dec{}3)\\
High-risk adenoma &
9\dec{}0 (4\dec{}9--13\dec{}7) &
8\dec{}4 (4\dec{}9--12\dec{}0) &
11\dec{}1 (5\dec{}7--17\dec{}1) &
10\dec{}4 (6\dec{}0--15\dec{}3)\\
Colorectal cancer &
0\dec{}271 (0\dec{}040--0\dec{}558) &
0\dec{}369 (0\dec{}059--0\dec{}766) &
0\dec{}300 (0\dec{}049--0\dec{}623) &
0\dec{}408 (0\dec{}063--0\dec{}851)\\
\bottomrule
\end{tabularx}
\begin{tablenotes}[flushleft]
\footnotesize
\item[] Values are simulation means (95\% HPD intervals). Adjusted population
detection corrects $d_j$ for FIT sensitivity and colonoscopy participation.
Modelled prevalence additionally corrects for colonoscopy sensitivity.
Uncertainty in colonoscopy participation was propagated.
\item[] FIT=faecal immunochemical test; HPD=highest posterior density.
\end{tablenotes}
\end{threeparttable}
\end{table}

\begin{figure}[p]
\centering
\includegraphics[width=\textwidth,height=0.78\textheight,keepaspectratio]
  {outputs/figure_publication_lancet_adherence_sampled_20261007/figure_adjusted_detection_prevalence_adherence_sampled.pdf}
\caption{\textbf{Adjusted population detection and modelled prevalence under
PRENEC-derived and literature-based FIT sensitivity scenarios.}
Panel A shows population detection after correction for FIT sensitivity and
colonoscopy participation; panel B additionally corrects for colonoscopy
sensitivity. Columns show low-risk adenoma, high-risk adenoma, and colorectal
cancer. The two rows identify the PRENEC-derived and literature-based FIT
sensitivity scenarios. Within each scenario, filled circles and solid
intervals denote Base A, open circles and dashed intervals denote Base B, and
Base A is shown above Base B.
Points are simulation means and horizontal lines are 95\% HPD intervals.
Uncertainty in colonoscopy participation was propagated in both panels.
FIT=faecal immunochemical test; HPD=highest posterior density.}
\label{fig:modelled}
\end{figure}

\subsection*{Results by sex}

In assignment A with PRENEC-derived FIT sensitivity, modelled prevalence was
higher among men than women for all three outcomes. Low-risk adenoma
prevalence was 36\dec{}5\% (95\% HPD interval 9\dec{}2--69\dec{}5) among men
and 22\dec{}1\% (7\dec{}6--38\dec{}7) among women. High-risk adenoma
prevalence was 16\dec{}4\% (5\dec{}6--28\dec{}7) among men and 9\dec{}3\%
(4\dec{}0--15\dec{}1) among women. Colorectal cancer prevalence was
0\dec{}531\% (0\dec{}032--1\dec{}249) among men and 0\dec{}215\%
(0\dec{}014--0\dec{}510) among women. Literature-based sensitivity produced
the same ordering. Uncertainty intervals were wide and overlapped between
sexes (table~\ref{tab:sexresults}).

\section*{Discussion}

\begin{center}
\fbox{\parbox{0.90\textwidth}{\centering
\textit{Placeholder. The Discussion will be drafted after author review of
the final model specification, results, and interpretation.}}}
\end{center}

\clearpage
\section*{Supplementary Methods}

\subsection*{Relationship with the Rogan--Gladen correction}

Let $Z_j=1$ denote completion of the positive sequential pathway: the
selected FIT was positive, colonoscopy was completed, and lesion category
$j$ was detected and histologically confirmed. Thus
$d_j=\Pr(Z_j=1)$. Defining

\begin{equation*}
\Se_{\mathrm{seq},j}=\Pr(Z_j=1\mid L_j=1)
\end{equation*}

and

\begin{equation*}
\Sp_{\mathrm{seq},j}=\Pr(Z_j=0\mid L_j=0),
\end{equation*}

the law of total probability gives

\begin{equation}
d_j=p_j\Se_{\mathrm{seq},j}
+(1-p_j)(1-\Sp_{\mathrm{seq},j}).
\label{eq:total}
\end{equation}

The Rogan--Gladen inversion is

\begin{equation}
p_j=
\frac{d_j+\Sp_{\mathrm{seq},j}-1}
{\Se_{\mathrm{seq},j}+\Sp_{\mathrm{seq},j}-1}.
\label{eq:rg}
\end{equation}

Under the conditional assumptions in the main Methods,

\begin{equation*}
\Se_{\mathrm{seq},j}
=\Se_{\mathrm{FIT},j}\,\Adh\,\Se_{\mathrm{COL},j}.
\end{equation*}

Because a positive final outcome required histological confirmation, the
sequence was treated as effectively specific:
$\Sp_{\mathrm{seq},j}\approx1$. Equation~\eqref{eq:rg} then becomes

\begin{equation*}
p_j=
\frac{d_j}
{\Se_{\mathrm{FIT},j}\,\Adh\,\Se_{\mathrm{COL},j}},
\qquad
r_j=
\frac{d_j}{\Se_{\mathrm{FIT},j}\,\Adh}.
\end{equation*}

The PRENEC estimators are therefore the effectively perfect-final-specificity
case of the Rogan--Gladen correction applied to the complete
FIT--colonoscopy--histology pathway. FIT specificity affects the number of
colonoscopies generated but cannot create a histologically confirmed adenoma
or cancer in the numerator.

\subsection*{Sex-specific model estimates}

\setcounter{table}{0}
\renewcommand{\thetable}{S\arabic{table}}

\begin{longtable}{@{}L{0.12\textwidth}L{0.14\textwidth}
L{0.16\textwidth}L{0.22\textwidth}L{0.22\textwidth}@{}}
\caption{Modelled prevalence in assignment A by sex and FIT sensitivity
scenario}\label{tab:sexresults}\\
\toprule
\textbf{Sex} & \textbf{Lesion} & \textbf{FIT scenario} &
\textbf{Mean, \%} & \textbf{95\% HPD interval, \%}\\
\midrule
\endfirsthead
\toprule
\textbf{Sex} & \textbf{Lesion} & \textbf{FIT scenario} &
\textbf{Mean, \%} & \textbf{95\% HPD interval, \%}\\
\midrule
\endhead
Women & Low risk & PRENEC & 22\dec{}10 & 7\dec{}59--38\dec{}67\\
Women & Low risk & Literature & 20\dec{}00 & 8\dec{}52--32\dec{}57\\
Women & High risk & PRENEC & 9\dec{}31 & 4\dec{}01--15\dec{}14\\
Women & High risk & Literature & 8\dec{}73 & 4\dec{}30--13\dec{}69\\
Women & Cancer & PRENEC & 0\dec{}215 & 0\dec{}014--0\dec{}510\\
Women & Cancer & Literature & 0\dec{}293 & 0\dec{}017--0\dec{}692\\
\addlinespace
Men & Low risk & PRENEC & 36\dec{}50 & 9\dec{}25--69\dec{}49\\
Men & Low risk & Literature & 33\dec{}02 & 9\dec{}43--59\dec{}78\\
Men & High risk & PRENEC & 16\dec{}37 & 5\dec{}55--28\dec{}74\\
Men & High risk & Literature & 15\dec{}34 & 5\dec{}95--25\dec{}79\\
Men & Cancer & PRENEC & 0\dec{}531 & 0\dec{}032--1\dec{}249\\
Men & Cancer & Literature & 0\dec{}721 & 0\dec{}053--1\dec{}695\\
\bottomrule
\end{longtable}

\clearpage
\begin{thebibliography}{99}
\small
\raggedright
\setlength{\itemsep}{1pt}
\setlength{\parskip}{0pt}

\bibitem{Arnold2017}
Arnold M, Sierra MS, Laversanne M, Soerjomataram I, Jemal A, Bray F.
Global patterns and trends in colorectal cancer incidence and mortality.
\textit{Gut}. 2017;66:683--691.
\href{https://doi.org/10.1136/gutjnl-2015-310912}{doi:10.1136/gutjnl-2015-310912}.

\bibitem{IARC2026}
Ervik M, Laversanne M, Lam F, Colombet M, Ferlay J, Filho AM, Bray F.
Global Cancer Observatory: Cancer Today (2024 estimates, version 1.0).
Lyon: International Agency for Research on Cancer; 2026.
\url{https://gco.iarc.who.int/today}.

\bibitem{Wang2023}
Wang S, Yang Z, Sha F, Qi X, He Z, Szeto CH, et al.
Prevalence of incidental colorectal cancer and polyps in autopsies of
different populations: a systematic review with meta-regression analysis.
\textit{Eur J Epidemiol}. 2023;38:939--955.
\href{https://doi.org/10.1007/s10654-023-01041-0}{doi:10.1007/s10654-023-01041-0}.

\bibitem{Heitman2009}
Heitman SJ, Ronksley PE, Hilsden RJ, Manns BJ, Rostom A, Hemmelgarn BR.
Prevalence of adenomas and colorectal cancer in average risk individuals:
a systematic review and meta-analysis.
\textit{Clin Gastroenterol Hepatol}. 2009;7:1272--1278.
\href{https://doi.org/10.1016/j.cgh.2009.05.032}{doi:10.1016/j.cgh.2009.05.032}.

\bibitem{Zhao2019}
Zhao S, Wang S, Pan P, Xia T, Chang X, Yang X, et al.
Magnitude, risk factors, and factors associated with adenoma miss rate of
tandem colonoscopy: a systematic review and meta-analysis.
\textit{Gastroenterology}. 2019;156:1661--1674.e11.
\href{https://doi.org/10.1053/j.gastro.2019.01.260}{doi:10.1053/j.gastro.2019.01.260}.

\bibitem{Montalvan2024}
Montalvan-Sanchez EE, Norwood DA, Dougherty M, et al.
Colorectal cancer screening programs in Latin America: a systematic review
and meta-analysis.
\textit{JAMA Netw Open}. 2024;7:e2354256.
\href{https://doi.org/10.1001/jamanetworkopen.2023.54256}{doi:10.1001/jamanetworkopen.2023.54256}.

\bibitem{Silva2011}
Silva MdeL, Santander R, Gobelet J, et al.
Plan de tamizaje de c\'ancer colorrectal ("Mes del Colon") en la
Cl\'inica Alemana de Santiago de Chile.
\textit{Acta Gastroenterol Latinoam}. 2011;41:10--16.
\href{https://pubmed.ncbi.nlm.nih.gov/21539063/}{PMID:21539063}.

\bibitem{Okada2016}
Okada T, Tanaka K, Kawachi H, Ito T, Nishikage T, Odagaki T, et al.
International collaboration between Japan and Chile to improve detection
rates in colorectal cancer screening.
\textit{Cancer}. 2016;122:71--77.
\href{https://doi.org/10.1002/cncr.29715}{doi:10.1002/cncr.29715}.

\bibitem{Paris2003}
The Paris endoscopic classification of superficial neoplastic lesions:
esophagus, stomach, and colon: November 30 to December 1, 2002.
\textit{Gastrointest Endosc}. 2003;58:S3--43.
\href{https://doi.org/10.1016/S0016-5107(03)02159-X}{doi:10.1016/S0016-5107(03)02159-X}.

\bibitem{Vienna2000}
Schlemper RJ, Riddell RH, Kato Y, Borchard F, Cooper HS, Dawsey SM, et al.
The Vienna classification of gastrointestinal epithelial neoplasia.
\textit{Gut}. 2000;47:251--255.
\href{https://doi.org/10.1136/gut.47.2.251}{doi:10.1136/gut.47.2.251}.

\bibitem{Ito2016}
Ito T, Kawachi H, Pe\~naloza P, Z\'arate A, Ponce A, Kobayashi M, et al.
Protocolo Estandarizado de Anatom\'ia Patol\'ogica de Lesiones Polipoideas
en Proyecto de Prevenci\'on de Neoplasias Colorrectales (PRENEC) en Chile.
\textit{Gastroenterol Latinoam}. 2016;27:37--46.
\href{https://doi.org/10.0716/gastrolat2016n100005}{doi:10.0716/gastrolat2016n100005}.

\bibitem{Fanshawe2024}
Fanshawe TR, Nicholson BD, Perera R, Oke JL.
A review of methods for the analysis of diagnostic tests performed in
sequence.
\textit{Diagn Progn Res}. 2024;8:8.
\href{https://doi.org/10.1186/s41512-024-00175-3}{doi:10.1186/s41512-024-00175-3}.

\bibitem{RoganGladen1978}
Rogan WJ, Gladen B.
Estimating prevalence from the results of a screening test.
\textit{Am J Epidemiol}. 1978;107:71--76.
\href{https://doi.org/10.1093/oxfordjournals.aje.a112510}{doi:10.1093/oxfordjournals.aje.a112510}.

\bibitem{Park2010}
Park DI, Ryu S, Kim YH, Lee SH, Lee CK, Eun CS, et al.
Comparison of guaiac-based and quantitative immunochemical fecal occult
blood testing in a population at average risk undergoing colorectal cancer
screening.
\textit{Am J Gastroenterol}. 2010;105:2017--2025.
\href{https://doi.org/10.1038/ajg.2010.179}{doi:10.1038/ajg.2010.179}.

\bibitem{Redwood2016}
Redwood DG, Asay ED, Blake ID, Sacco PE, Christensen CM, Sacco FD, et al.
Stool DNA testing for screening detection of colorectal neoplasia in Alaska
Native people.
\textit{Mayo Clin Proc}. 2016;91:61--70.
\href{https://doi.org/10.1016/j.mayocp.2015.10.008}{doi:10.1016/j.mayocp.2015.10.008}.

\bibitem{Chang2017}
Chang LC, Shun CT, Hsu WF, Tu CH, Tsai PY, Lin BR, et al.
Fecal immunochemical test detects sessile serrated adenomas and polyps with
a low level of sensitivity.
\textit{Clin Gastroenterol Hepatol}. 2017;15:872--879.e1.
\href{https://doi.org/10.1016/j.cgh.2016.07.029}{doi:10.1016/j.cgh.2016.07.029}.

\bibitem{Imperiale2014}
Imperiale TF, Ransohoff DF, Itzkowitz SH, Levin TR, Lavin P, Lidgard GP,
et al.
Multitarget stool DNA testing for colorectal-cancer screening.
\textit{N Engl J Med}. 2014;370:1287--1297.
\href{https://doi.org/10.1056/NEJMoa1311194}{doi:10.1056/NEJMoa1311194}.

\bibitem{Lin2021}
Lin JS, Perdue LA, Henrikson NB, Bean SI, Blasi PR.
Screening for colorectal cancer: updated evidence report and systematic
review for the US Preventive Services Task Force.
\textit{JAMA}. 2021;325:1978--1998.
\href{https://doi.org/10.1001/jama.2021.4417}{doi:10.1001/jama.2021.4417}.

\bibitem{Martin2013}
Mart\'in L\'opez JE, M\'arquez Pel\'aez S, Adam Blanco D,
Navarro Caballero JA, Rodr\'iguez L\'opez R, Beltr\'an Calvo C.
\textit{Eficacia, seguridad y eficiencia de la colonograf\'ia por
tomograf\'ia computerizada frente a la colonoscopia como prueba de cribado
del c\'ancer colorrectal: revisi\'on sistem\'atica y metaan\'alisis}.
Sevilla: Agencia de Evaluaci\'on de Tecnolog\'ias Sanitarias de
Andaluc\'ia; 2013. Informe AETSA 2011/01.

\bibitem{vanRijn2006}
van Rijn JC, Reitsma JB, Stoker J, Bossuyt PM, van Deventer SJ, Dekker E.
Polyp miss rate determined by tandem colonoscopy: a systematic review.
\textit{Am J Gastroenterol}. 2006;101:343--350.
\href{https://doi.org/10.1111/j.1572-0241.2006.00390.x}{doi:10.1111/j.1572-0241.2006.00390.x}.

\bibitem{Zhao2023}
Zhao S, Song Y, Wang S, Wang R, Feng Z, Gong A, et al.
Reduced adenoma miss rate with 9-minute vs 6-minute withdrawal times for
screening colonoscopy: a multicenter randomized tandem trial.
\textit{Am J Gastroenterol}. 2023;118:802--811.
\href{https://doi.org/10.14309/ajg.0000000000002055}{doi:10.14309/ajg.0000000000002055}.

\bibitem{STROBE2007}
von Elm E, Altman DG, Egger M, Pocock SJ, G\o tzsche PC,
Vandenbroucke JP.
The Strengthening the Reporting of Observational Studies in Epidemiology
(STROBE) statement.
\textit{Lancet}. 2007;370:1453--1457.
\href{https://doi.org/10.1016/S0140-6736(07)61602-X}{doi:10.1016/S0140-6736(07)61602-X}.

\bibitem{RECORD2015}
Benchimol EI, Smeeth L, Guttmann A, et al.
The REporting of studies Conducted using Observational Routinely-collected
health Data (RECORD) statement.
\textit{PLoS Med}. 2015;12:e1001885.
\href{https://doi.org/10.1371/journal.pmed.1001885}{doi:10.1371/journal.pmed.1001885}.

\bibitem{GATHER2016}
Stevens GA, Alkema L, Black RE, et al.
Guidelines for Accurate and Transparent Health Estimates Reporting: the
GATHER statement.
\textit{PLoS Med}. 2016;13:e1002056.
\href{https://doi.org/10.1371/journal.pmed.1002056}{doi:10.1371/journal.pmed.1002056}.

\end{thebibliography}

\end{document}
